\subsection{Every weight-preserving Frobenius has the same completion defect}

\begin{proposition}[A-WF02: Proved; all weight-preserving Dieudonn\'e Frobenii in the stated coefficient setting]
\label{AWF02}
Fix a prime $p$ and let $R$ be $\mathbb Z$,
$\mathbb Z_{(p)}$, or $\mathbb Z_p$.
Put $S=\Gamma_R(T)$, with coefficient weight one
and form degree zero. Let $A$ be a nonnegative,
weight-graded, simplicial strict PD differential
algebra whose terms are $p$-torsionfree. Suppose
its actual closed-zero-form algebra is $S$ in
every simplex and every weight. Suppose also that
each simplex has a weight-preserving vertex
retraction $\epsilon_n:A_n^0\to S$ of the
coefficient inclusion; in particular this holds
when $A_0=S$ and the retractions are vertex maps.

Let $F$ be any simplicially natural, weight-preserving,
unital graded-ring endomorphism that makes each
$A_n$ a Dieudonn\'e algebra for the prime $p$.
No PD compatibility of $F$ is assumed.
There is a single $c\in pR$ such that, in every
simplex,
\[
                    F(T^{[w]})=c^w T^{[w]}
                    \quad(w\geq0).
 \tag{WF.9}
\]
The actual canonical composite of Dieudonn\'e
algebras
\[
 A\xrightarrow{\ j\ }\operatorname{Sat}_F(A)
     \longrightarrow\mathcal W\operatorname{Sat}_F(A)
 \tag{WF.10}
\]
sends every positive-weight coefficient to zero.
Consequently this completion with its induced
coefficient map cannot give an admissible theory
with the specified coefficient comparison.
This statement covers every $F$ satisfying the
displayed weight-preserving assumptions.
\end{proposition}

\begin{proof}
We use precisely the opened primary constructions
of Bhatt--Lurie--Mathew
\cite{Root_BhattLurieMathew2020}: Definition~3.1.2
for Dieudonn\'e algebras; Definition~2.2.1,
Remark~2.2.3 and Proposition~2.2.4 in the v3 source,
for saturated
complexes and Verschiebung; and Construction~2.5.1
for completion. The exact source lines used here
are respectively 1404--1413, 615--645, and 860--869.
In particular, on a saturated complex $F$ is
injective, $FV=VF=p$, and the completion is the
inverse limit of the actual quotients by
$\operatorname{im}(V^r)+\operatorname{im}(dV^r)$.

The equality $dF=pFd$ makes the image of a closed
zero-form closed. Weight preservation and the
given closed-zero-form identification therefore
give $F(T)=c_n T$ for a scalar $c_n\in R$ in
simplex $n$. Applying a simplicial vertex map
and using naturality of $F$ shows $c_n=c_0$:
both coefficient copies inject and the vertex
map is the identity on their marked $T$.
We write this common scalar as $c$.

The restriction of $F$ to the closed weight-zero
coefficient ring is the identity on $R$. For
$R=\mathbb Z$ this follows from its unit; for
$\mathbb Z_{(p)}$ it also follows on every inverted
integer. For $R=\mathbb Z_p$, the restriction
maps $R$ to $R$, and for every $a\in R$ and
integer representative $a_r$ modulo $p^r$ one
has $F(a)\equiv a_r\equiv a\pmod{p^r}$.
The intersection of the ideals $p^r\mathbb Z_p$
is zero, giving the same identity.

The degree-zero Frobenius congruence gives
$F(T)\equiv T^p\pmod{pA_n^0}$. Since
$T^p=p!T^{[p]}$, its right-hand side is zero
modulo $p$. Apply the actual retraction
$\epsilon_n$ to conclude $cT\in pS$.
The coefficient basis of $S$ then gives $c\in pR$.
For every $w$, the element $F(T^{[w]})$ is closed
and of weight $w$, hence is $c_w T^{[w]}$ for
some $c_w\in R$. Applying the multiplicative $F$
to $T^w=w!T^{[w]}$ gives
\[
                     w!c_w T^{[w]}
                          =w!c^wT^{[w]}.
 \tag{WF.11}
\]
Cancellation takes place in the torsionfree
coefficient ring $S$, where multiplication by
$w!$ is injective for all three rings $R$.
This proves (WF.9), without using a divided-power
identity for $F$ itself.

Put $M=\operatorname{Sat}_F(A)$ and let
$a=j(T^{[w]})\in M_n^0$, with $w>0$.
For every integer $r\geq1$ the scalar
\[
            \frac{c^{rw}}{p^r}
             =p^{r(w-1)}(c/p)^{rw}
 \tag{WF.12}
\]
belongs to $R$. Define the actual element
$b_r=j((c^{rw}/p^r)T^{[w]})\in M_n^0$.
The map $j$ respects $F$, and (WF.9) with the
identity on $R$ gives $F^r(a)=c^{rw}a$.
The known saturated identities consequently give
\[
 F^r(V^r b_r)=p^r b_r=c^{rw}a=F^r(a).
 \tag{WF.13}
\]
Injectivity of $F^r$ proves the exact equality
\[
                              a=V^r b_r.
 \tag{WF.14}
\]
The argument includes $c=0$: in that case it
also proves that this coefficient already has
zero image in the saturation.

Thus $a$ has zero image in each quotient defining
$\mathcal W(M)$. Its image in their inverse
limit is zero, proving (WF.10). The calculation
uses the actual canonical maps after regarding
the given PD algebra as a Dieudonn\'e algebra;
it requires no assertion of an additional PD
structure on a completed object. Any proposed
admissible PD theory with this very induced
coefficient map would have $T^{[w]}$ in that
map's kernel for every $w>0$, whereas $T^{[w]}$
is nonzero in $S$. No eventual weight cutoff
removes that explicit kernel.
\end{proof}
